\begin{problem}[Warmup]{(a)}
    Return to Exercise a in the previous chapter. In each case, determine a basis for the simple extension obtained by adjoining the element to the field in question.
\end{problem}
\unsolved

\begin{problem}[Warmup]{(b)}
    Express the following field elements as a linear combination of the basis elements you determined in Exercise a.
    \begin{enumerate}
        \item[(a)] $\dfrac{1}{2 + \sqrt{2}} \in \mathbb{Q}(1 + \sqrt{2})$.
        \item[(b)] $\dfrac{\sqrt[3]{2}}{1 + 2\sqrt[3]{2}} \in \mathbb{Q}(\sqrt[3]{2})$.
    \end{enumerate}
\end{problem}
\unsolved

\begin{problem}[Warmup]{(c)}
    Explain why the following statements are true, or else give a counterexample. In each case $F$ is a subfield of the field $E$, and $\alpha \in E \setminus F$:
    \begin{enumerate}
        \item[(a)] $F(\alpha)$ is finite dimensional over $F$.
        \item[(b)] Let $\alpha$ be algebraic over $F$. Then $F(\alpha)$ is finite dimensional over $F$.
        \item[(c)] Every element of $F(\alpha)$ is algebraic over $F$.
        \item[(d)] Every element of $F(\alpha)$ is algebraic over $E$.
        \item[(e)] Suppose that $\alpha$ is algebraic over $F$. Then every element of $F(\alpha)$ is algebraic over $F$.
    \end{enumerate}
\end{problem}
\unsolved

\begin{problem}[Warmup]{(d)}
    Explain why the minimal polynomial is unique.
\end{problem}
\unsolved

\begin{problem}{1}
    \begin{enumerate}
        \item[(a)] Compute the addition and multiplication tables for the field
        $\mathbb{Z}_2(\alpha)$ described in Example 43.7. Use your table to determine
        explicitly the multiplicative inverse of each non-zero element of this field. (This
        exercise is essentially a repeat of Exercise 8.12.)
        \item[(b)] In Exercise 17.18, you considered the Frobenius isomorphism, defined on
        any finite field. Compute this homomorphism explicitly in this case.
    \end{enumerate}
\end{problem}
\unsolved

\begin{problem}{2}
    Consider the polynomial
    \[f = 1 + x^2 + x^3 \in \mathbb{Z}_2(\alpha)[x],\]
    where $\mathbb{Z}_2(\alpha)$ is the field considered in the previous exercise.
    \begin{enumerate}
        \item[(a)] Use the Root Theorem to show that this polynomial is irreducible.
        \item[(b)] By Kronecker's Theorem 42.1 we can construct an extension field of
        $\mathbb{Z}_2(\alpha)$ so that $f$ has a root $\beta$. How many elements does this
        field have?
        \item[(c)] If you're not doing anything this weekend, construct a multiplication
        table for this field.
    \end{enumerate}
\end{problem}
\unsolved

\begin{problem}{3}
    In this problem you will prove that
    \[f = x^6 - 6x^4 - 4x^3 + 12x^2 - 24x - 4\]
    is the minimal polynomial for $\alpha = \sqrt{2} + \sqrt[3]{2}$, as stated in
    Example 43.6. (We will return to this example in Example 44.3.)
    \begin{enumerate}
        \item[(a)] Check that $f(\alpha) = 0$.
        \item[(b)] Show that $f$ is irreducible in $\mathbb{Q}[x]$, by using Eisenstein's
        Criterion 5.7, after the linear change of variables $x = y + 2$ (see
        Exercise 5.14).
        \item[(c)] Show that $f$ can be factored into irreducibles in $\mathbb{R}[x]$ as
        follows:
        \begin{align*}
            &\left(x - \sqrt{2} - \sqrt[3]{2}\right)
            \left(x^2 - (2\sqrt{2} + \sqrt[3]{2})x + (2 + \sqrt{2}\sqrt[3]{2} + \sqrt[3]{4})\right) \\
            &\left(x + \sqrt{2} - \sqrt[3]{2}\right)
            \left(x^2 + (2\sqrt{2} - \sqrt[3]{2})x + (2 - \sqrt{2}\sqrt[3]{2} + \sqrt[3]{4})\right).
        \end{align*}
        \item[(d)] How would you factor $f$ into irreducibles in $\mathbb{C}[x]$? (You need
        not feel obliged to actually carry this out!)
    \end{enumerate}
\end{problem}
\unsolved

\begin{problem}{4}
    Determine the minimal polynomial of $\sqrt{2} + \sqrt{3}$ over the following three
    fields: $\mathbb{Q}$, $\mathbb{Q}(\sqrt{2})$, $\mathbb{R}$.
\end{problem}
\unsolved

\begin{problem}{5}
    Let $\alpha = \sqrt{2} + \sqrt[3]{2}$, and $f$ be its minimal polynomial, as in
    Example 43.6, and in Exercise 3. Show explicitly that the following elements of
    $\mathbb{Q}(\alpha)$ are linear combinations of
    \[1, \alpha, \alpha^2, \ldots, \alpha^5.\]
    \begin{enumerate}
        \item[(a)] $\dfrac{1}{\alpha + 1}$.
        \item[(b)] $\dfrac{\alpha}{1 + \alpha^2}$.
    \end{enumerate}
\end{problem}
\unsolved

\begin{problem}{6}
    Let
    \[K = \left\{\frac{f(\pi)}{g(\pi)} \in \mathbb{R} : f, g \in \mathbb{Q}[x],\ g \neq 0\right\},\]
    as given in Example 43.10.
    \begin{enumerate}
        \item[(a)] Why do the quotients defined above always make sense? (That is, why are
        the denominators always non-zero?)
        \item[(b)] Show that $K$ is a subfield of $\mathbb{R}$.
        \item[(c)] Argue that $K = \mathbb{Q}(\pi)$.
        \item[(d)] Are the elements in $K$ uniquely represented, as we've presented them in
        the definition of $K$?
    \end{enumerate}
\end{problem}
\unsolved

\begin{problem}{7}
    Let $F$ be a subfield of the field $E$, and suppose that $\alpha \in E$ is
    transcendental over $F$. Using Exercise 6 as a model, formulate an explicit description
    of the members of the simple extension $F(\alpha)$, and prove that your formulation
    works.
\end{problem}
\unsolved

\begin{problem}{8}
    Let $F$ be a subfield of the field $E$, and suppose that $\alpha \in E$ is 
        transcendental over $F$.
    \begin{enumerate}
        \item Show that 
        \[1, \alpha, \alpha^2, \alpha^3, \cdots\]
        is a linearly independent set of the vector space $F(\alpha)$ (over $F$).
        \item Why does the previous part mean that the dimension of $F(\alpha)$ over $F$ is infinite.
        \item Show that the set in the first part is \emph{not} a basis for $F(\alpha)$ over $F$.
    \end{enumerate}
\end{problem}
\solved

\begin{proof}
    Suppose for contradiction that $1, \alpha, \alpha^2, \ldots$ are linearly
    dependent over $F$. 
    There exist $n \ge 0$ and $b_0, \ldots, b_n \in F$, not all zero, with
    \[
        b_0 + b_1 \alpha + \cdots + b_n \alpha^n = 0.
    \]
    Let $f(x) = b_0 + b_1 x + \cdots + b_n x^n \in F[x]$. Since some $b_i \ne 0$,
    $f$ is a nonzero polynomial, and $f(\alpha) = 0$. Thus
    $\alpha$ is algebraic over $F$, contradicting that $\alpha$ is transcendental.
\end{proof}

\begin{proof}
    If $F(\alpha)$ had finite dimension $n$ over $F$, then every linearly
    independent subset of $F(\alpha)$ would have at most $n$ elements. But
    $1, \alpha, \ldots, \alpha^n$ are $n + 1$ linearly independent elements
    by the first part, a contradiction.
\end{proof}

\begin{proof}
    Since $0$ is algebraic over $F$, $\alpha \ne 0$, so $\alpha^{-1} \in F(\alpha)$.
    Suppose for contradiction that $\alpha^{-1} = b_0 + b_1 \alpha + \cdots + b_n \alpha^n$
    for some $b_0, \ldots, b_n \in F$. Multiplying by $\alpha$ gives
    \[
        b_n \alpha^{n+1} + \cdots + b_0 \alpha - 1 = 0,
    \]
    so $\alpha$ is a root of $g(x) = b_n x^{n+1} + \cdots + b_0 x - 1 \in F[x]$.
    The constant term of $g$ is $-1$, so $g$ is nonzero, and $\alpha$ is algebraic
    over $F$, a contradiction. Thus $\alpha^{-1}$ is not in the span of
    $1, \alpha, \alpha^2, \ldots$, so this set is not a basis of $F(\alpha)$.
\end{proof}
\begin{problem}{9}
    In this exercise we perform some additional verifications related to Example 43.9.
    Show that $\sqrt[3]{2}, \sqrt[3]{2}\zeta \notin \mathbb{Q}(\sqrt[3]{2}\zeta^2)$ and
    $\sqrt[3]{2}\zeta^2 \notin \mathbb{Q}(\sqrt[3]{2}\zeta)$. Thus, in each case the field
    isomorphisms between the three simple extensions of Example 43.9 are very simply
    described. How?
\end{problem}
\unsolved
